% (C) 2026 Ralf Stephan, in collaboration with Claude Code.
% Released under CC0 1.0 Universal (public-domain dedication).
% See https://creativecommons.org/publicdomain/zero/1.0/

\documentclass[11pt]{article}
\usepackage[T1]{fontenc}
\usepackage[utf8]{inputenc}
\usepackage{lmodern}
\usepackage[margin=2.5cm]{geometry}
\usepackage{amsmath,amssymb,amsthm}
\usepackage[hidelinks]{hyperref}

% Lean names and paths: verbatim, breakable at . _ /
\DeclareUrlCommand\lean{\urlstyle{tt}}
\def\UrlBreaks{\do\.\do\_\do\/}
\emergencystretch=6em

\newtheorem{theorem}{Theorem}[section]
\newtheorem{lemma}[theorem]{Lemma}
\newtheorem{corollary}[theorem]{Corollary}
\newtheorem{proposition}[theorem]{Proposition}
\theoremstyle{definition}
\newtheorem{definition}[theorem]{Definition}
\theoremstyle{remark}
\newtheorem{remark}[theorem]{Remark}

\title{Further results on the brief\\ \emph{Schr\"odinger equation is equivalent to information
  conservation}}
\author{Ralf Stephan}
\date{16 November 2026}

\begin{document}
\maketitle

\begin{abstract}
  This note collects results that go beyond the technical brief \emph{Schr\"odinger equation is
  equivalent to information conservation}~\cite{brief}. They are not part of the brief. Each
  section names the Lean declarations that check it; the Lean files are in the brief's
  directory, under \lean{UnitaryEntropy/FurtherResults/}.
\end{abstract}

\section{The classical counterpart}
\label{sec:classical}

The brief's conclusion says that its result ``parallels the classical case where Hamiltonian
evolution is equivalent to the preservation of Shannon entropy (Liouville's theorem) and of the
definitional independence of degrees of freedom''. This section states and proves that classical
result, along the lines of the argument in~\cite[Ch.~\emph{Classical mechanics}, Sec.~\emph{Multiple
degrees of freedom}]{book}. The proof runs parallel to the brief's: entropy fixes a determinant,
independence fixes the form $\omega$ up to a factor $\pm1$, continuity in time excludes $-1$, and
an infinitesimal argument gives the equation of motion.

\subsection{Statement}

The phase space of $n$ degrees of freedom is $\mathbb R^{2n}$ with coordinates
$\xi = (q_1, \dots, q_n, p_1, \dots, p_n)$ and the symplectic form
\[
  \omega(u, v) = \sum_i \bigl(u_{q_i} v_{p_i} - u_{p_i} v_{q_i}\bigr).
\]
An \emph{evolution} is a family of maps $\varphi_{t_0\to t_1}\colon \mathbb R^{2n}\to\mathbb R^{2n}$,
for all times $t_0, t_1$, that is deterministic and reversible,
\[
  \varphi_{t\to t} = \mathrm{id}, \qquad \varphi_{t_1\to t_2}\circ\varphi_{t_0\to t_1} =
  \varphi_{t_0\to t_2},
\]
so that $\varphi_{t_1\to t_0}$ reconstructs the past, and differentiable: $(t, \xi) \mapsto
\varphi_{t_0\to t}(\xi)$ is of class $C^2$ for every $t_0$. Its \emph{displacement field} is
$S(t, \xi) = \partial_s \varphi_{t\to s}(\xi)|_{s=t}$, and every orbit moves with it:
$\partial_t \varphi_{t_0\to t}(\xi) = S(t, \varphi_{t_0\to t}(\xi))$. Write $D\varphi$ for the
derivative in $\xi$.

\begin{description}
\item[HM.] Hamilton's equations: there is $H(t, \xi)$, differentiable in $\xi$, with
  $dH_t = \omega(S(t, \cdot), \cdot)$, that is, $\dot q_i = \partial H/\partial p_i$ and
  $\dot p_i = -\partial H/\partial q_i$ along every orbit.
\item[DR-INFO.] The evolution preserves the Shannon entropy
  $I[\rho] = -\int \rho\log\rho\,d\xi$ of every probability density $\rho$. The evolved density
  is $\rho_{t_1}(\varphi(\xi)) = \rho(\xi)/|\det D\varphi(\xi)|$, $\varphi = \varphi_{t_0\to t_1}$.
\item[DR-IND.] The evolution preserves the independence of degrees of freedom: if
  $\omega(u, v) = 0$, then $\omega(D\varphi\,u, D\varphi\,v) = 0$.
\end{description}

In DR-IND, two degrees of freedom are independent exactly when $\omega$ vanishes across them,
and $\omega$ measures the area within each~\cite{book}. Asking that $\omega$-orthogonal
displacements stay $\omega$-orthogonal is the infinitesimal form. For $n \ge 2$ it is the same as
asking that independent degrees of freedom stay independent: two linearly independent
$\omega$-orthogonal vectors lie in two independent degrees of freedom (extend them to a
symplectic basis), and for dependent ones there is nothing to preserve. This equivalence is not
formalized.

\begin{theorem}\label{thm:classical}
  For an evolution as above, HM $\Leftrightarrow$ DR-INFO $+$ DR-IND. Both are equivalent to the
  evolution being symplectic, $\omega(D\varphi\,u, D\varphi\,v) = \omega(u, v)$, and then
  $\det D\varphi = 1$ (Liouville).
\end{theorem}

\begin{corollary}\label{cor:one}
  For one degree of freedom, HM $\Leftrightarrow$ DR-INFO.
\end{corollary}

\begin{proposition}\label{prop:two}
  For two degrees of freedom neither condition suffices alone. The squeeze
  $\varphi_{t_0\to t_1}(q_1, q_2, p_1, p_2) = (e^{t_1-t_0}q_1, e^{t_0-t_1}q_2, p_1, p_2)$ satisfies
  DR-INFO but not DR-IND, and the dilation $\varphi_{t_0\to t_1}(\xi) = e^{t_1-t_0}\xi$ satisfies
  DR-IND but not DR-INFO. Neither is Hamiltonian.
\end{proposition}

The squeeze breaks the independence of $q_1 + q_2$ and $p_1 - p_2$. Like the book's example of
linear drag, the dilation multiplies the area within every degree of freedom by the same
factor.

\subsection{Proof}

\begin{lemma}[Entropy]\label{lem:entropy}
  Let $\varphi$ be injective and $C^1$ with $\det D\varphi \ne 0$. Then $\varphi$ preserves the
  Shannon entropy of every probability density iff $|\det D\varphi| = 1$ everywhere.
\end{lemma}

\begin{proof}
  By the change of variables formula, $I[\rho_{t_1}] = I[\rho] + \int\rho\log|\det D\varphi|$.
  If $|\det D\varphi| = 1$ the last term vanishes. If not, $g = \log|\det D\varphi|$ is continuous
  and $g(\xi_0) \ne 0$ somewhere, so $g(\xi_0)\,g > g(\xi_0)^2/2$ on a ball $B$ around $\xi_0$.
  The uniform density on $B$ then has $\int\rho g \ne 0$.
\end{proof}

\begin{lemma}[Same orthogonality]\label{lem:ortho}
  Let $B$ and $B'$ be bilinear forms on a vector space over a field, with $B(u, \cdot) \ne 0$ for
  $u \ne 0$. If $B'(u, v) = 0$ whenever $B(u, v) = 0$, then $B' = cB$ for a scalar $c$.
\end{lemma}

\begin{proof}
  For $u \ne 0$ the functional $B'(u, \cdot)$ vanishes on the kernel of $B(u, \cdot)$, so
  $B'(u, \cdot) = c_u B(u, \cdot)$; likewise $B'(\cdot, v) = d_v B(\cdot, v)$ when
  $B(\cdot, v) \ne 0$. If $B(u, v) \ne 0$, comparing the two at $(u, v)$ gives $c_u = d_v$. For
  $u, u' \ne 0$ there is a $v$ with $B(u, v) \ne 0$ and $B(u', v) \ne 0$ (a vector space is not
  the union of two proper subspaces), so $c_u = d_v = c_{u'}$.
\end{proof}

\begin{lemma}[Determinant]\label{lem:det}
  If a linear map $A$ of $\mathbb R^{2n}$ satisfies $\omega(Au, Av) = c\,\omega(u, v)$, then
  $(\det A)^2 = c^{2n}$.
\end{lemma}

\begin{proof}
  In matrices, $A^{\mathsf T}JA = cJ$; take determinants and cancel $\det J \ne 0$.
\end{proof}

\begin{lemma}[Variational equation]\label{lem:variational}
  $DS(t, \xi) = \partial_s D\varphi_{t\to s}(\xi)|_{s=t}$, and
  $\partial_t D\varphi_{t_0\to t}(\xi) = DS(t, \varphi_{t_0\to t}(\xi))\,D\varphi_{t_0\to t}(\xi)$.
\end{lemma}

\begin{proof}
  The first equation exchanges the derivatives in $s$ and $\xi$, which is allowed for a $C^2$
  function. The second follows from the first, applied at time $t$ to
  $D\varphi_{t_0\to s} = D\varphi_{t\to s}(\varphi_{t_0\to t}(\xi))\,D\varphi_{t_0\to t}(\xi)$.
\end{proof}

\begin{proof}[Proof of Theorem~\ref{thm:classical}]
  DR-INFO $+$ DR-IND $\Rightarrow$ symplectic. Fix $t_0$ and $\xi$, and let
  $A_t = D\varphi_{t_0\to t}(\xi)$; it is invertible, since $\varphi_{t\to t_0}$ inverts
  $\varphi_{t_0\to t}$. By Lemma~\ref{lem:entropy}, $|\det A_t| = 1$. By Lemma~\ref{lem:ortho}
  with $B = \omega$ and $B'(u, v) = \omega(A_t u, A_t v)$, $\omega(A_t u, A_t v) =
  \kappa(t)\,\omega(u, v)$, and by Lemma~\ref{lem:det} $\kappa(t)^{2n} = 1$, so $\kappa(t) = \pm1$.
  This is the analogue of Wigner's theorem in the brief's Step~2, where the map is unitary or
  anti-unitary. Here $\kappa(t) = \omega(A_t e_{q_1}, A_t e_{p_1})$ is continuous in $t$, by
  Lemma~\ref{lem:variational}, and $\kappa(t_0) = 1$. By the intermediate value theorem
  $\kappa \equiv 1$, as joint continuity excludes the anti-unitary case in the brief.

  Symplectic $\Rightarrow$ HM. Differentiating $\omega(D\varphi_{t\to s}a, D\varphi_{t\to s}b) =
  \omega(a, b)$ at $s = t$, by Lemma~\ref{lem:variational},
  $\omega(DS\,a, b) + \omega(a, DS\,b) = 0$, that is, $\omega(DS\,a, b) = \omega(DS\,b, a)$. So
  the 1-form $\alpha_t = \omega(S(t, \cdot), \cdot)$ has a symmetric derivative: it is closed.
  On $\mathbb R^{2n}$, which is convex, the Poincar\'e lemma gives $H(t, \cdot)$ with
  $dH_t = \alpha_t$. This is the analogue of the brief's Step~3.

  HM $\Rightarrow$ symplectic. The second derivative of $H(t, \cdot)$ is symmetric, so
  $\omega(DS\,a, b) = \omega(DS\,b, a)$ as above. By Lemma~\ref{lem:variational}, the derivative
  in $t$ of $\omega(A_t u, A_t v)$ is $\omega(DS\,A_tu, A_tv) + \omega(A_tu, DS\,A_tv) = 0$, so it
  stays $\omega(u, v)$.

  Symplectic $\Rightarrow$ DR-INFO $+$ DR-IND. DR-IND is immediate. The matrix of $D\varphi$ is
  symplectic, so $\det D\varphi = 1$, and Lemma~\ref{lem:entropy} gives DR-INFO.
\end{proof}

\begin{proof}[Proof of Corollary~\ref{cor:one}]
  For $n = 1$, writing $u = u_q e_q + u_p e_p$ and expanding,
  $\omega(Au, Av) = \omega(Ae_q, Ae_p)\,\omega(u, v)$ for every linear $A$. So DR-IND always holds.
\end{proof}

\begin{proof}[Proof of Proposition~\ref{prop:two}]
  The squeeze is diagonal with $\det D\varphi = e^{s}e^{-s} = 1$, $s = t_1 - t_0$, so it satisfies
  DR-INFO by Lemma~\ref{lem:entropy}. For $u = e_{q_1} + e_{q_2}$ and $v = e_{p_1} - e_{p_2}$,
  $\omega(u, v) = 0$ but $\omega(D\varphi\,u, D\varphi\,v) = e^{s} - e^{-s} \ne 0$ for $s = 1$.
  The dilation has $\omega(D\varphi\,u, D\varphi\,v) = e^{2s}\omega(u, v)$, so it satisfies
  DR-IND but is not symplectic. By Theorem~\ref{thm:classical}, neither is Hamiltonian, and the
  dilation does not satisfy DR-INFO.
\end{proof}

\begin{remark}
  The differentiability of the evolution is assumed on both sides of Theorem~\ref{thm:classical}.
  In the brief, EVDIFF is part of the conclusion of the direction DR-SCEQ $\Rightarrow$ DR-INFO;
  classically it would follow from the smooth dependence of the solutions of Hamilton's equations
  on the initial values, for a smooth $H$, which Mathlib does not have. The Hamiltonian
  $H(t, \cdot)$ is constructed separately for each $t$, with no regularity in $t$, as the
  brief's $H(t)$.
\end{remark}

\subsection{Lean}

Module \lean{UnitaryEntropy.FurtherResults.ClassicalCounterpart}, namespace
\lean{UnitaryEntropy.FurtherResults}. It does not import the brief's files. The phase space is
\lean{PhaseSpace}~$\iota$ $= \iota \oplus \iota \to \mathbb R$ for a finite type $\iota$ of degrees
of freedom, with $q_i$ at \lean{inl}~$i$ and $p_i$ at \lean{inr}~$i$.

\begin{itemize}
  \item \lean{symplecticForm}: $\omega$, with matrix $-J$ for Mathlib's \lean{Matrix.J}.
  \item \lean{IsFlow}, \lean{displacement}: evolutions and the displacement field $S$;
    \lean{IsFlow.hasDerivAt_orbit}: orbits move with $S$.
  \item \lean{IsHamiltonian}, \lean{PreservesEntropy}, \lean{PreservesIndependence},
    \lean{IsSymplectic}: the conditions HM, DR-INFO, DR-IND and the symplectic condition;
    \lean{evolveDensity}: the evolved density. \lean{IsFlow.hamilton_equations}: HM gives
    Hamilton's equations along every orbit.
  \item \lean{MeasureTheory.integral_negMulLog_pushDensity_of_abs_det_eq_one},
    \lean{MeasureTheory.abs_det_eq_one_of_integral_negMulLog_pushDensity}:
    Lemma~\ref{lem:entropy}, in \lean{ForMathlib/InformationTheory/EntropyPreserving.lean}.
  \item \lean{LinearMap.exists_eq_smul_of_forall_isOrtho_imp}: Lemma~\ref{lem:ortho}, in
    \lean{ForMathlib/LinearAlgebra/BilinearMap/OrthogonalityPreserving.lean}.
  \item \lean{det_sq_eq_of_symplecticForm}: Lemma~\ref{lem:det}.
  \item \lean{hasDerivAt_fderiv_partial}, \lean{hasFDerivAt_deriv_partial},
    \lean{IsFlow.hasDerivAt_fderiv}: Lemma~\ref{lem:variational}.
  \item \lean{IsFlow.isSymplectic_of_preservesEntropy_of_preservesIndependence},
    \lean{IsFlow.isHamiltonian_of_isSymplectic}, \lean{IsFlow.isSymplectic_of_isHamiltonian},
    \lean{IsSymplectic.det_eq_one}, \lean{IsFlow.isSymplectic_iff},
    \lean{IsFlow.isHamiltonian_iff_isSymplectic}, \lean{IsFlow.isHamiltonian_iff}:
    Theorem~\ref{thm:classical}. The Poincar\'e lemma is Mathlib's
    \lean{Convex.exists_forall_hasFDerivAt_of_fderiv_symmetric}, and $\det = 1$ for symplectic
    matrices is Mathlib's \lean{SymplecticGroup.det_eq_one}.
  \item \lean{symplecticForm_map_of_unique}, \lean{IsFlow.isHamiltonian_iff_preservesEntropy}:
    Corollary~\ref{cor:one}.
  \item \lean{squeezeFlow_counterexample}, \lean{dilationFlow_counterexample}:
    Proposition~\ref{prop:two}. The dilation is stated for any nonempty set of degrees of
    freedom.
\end{itemize}

\section{Entropy conservation as determinism and reversibility}
\label{sec:detrev}

The brief reads DR-INFO as ``an information theoretic way to characterize determinism and
reversibility'', but defines neither notion. This section gives a precise definition and proves
it equivalent to DR-INFO, for a single evolution map and without continuity in time.

\subsection{Statement}

The setting is the brief's: a finite-dimensional complex Hilbert space $\mathcal H$ (QST:H), the
density operators $\mathcal D(\mathcal H)$, the pure states
$\mathcal P(\mathcal H) = \{|\psi\rangle\langle\psi| : \|\psi\| = 1\}$, and a map
$\Phi\colon\mathcal D(\mathcal H)\to\mathcal D(\mathcal H)$ that preserves equal mixtures,
$\Phi(\frac12\rho + \frac12\sigma) = \frac12\Phi(\rho) + \frac12\Phi(\sigma)$ ([B]EVMIX). The
entropy is the von Neumann entropy $S(\rho) = -\operatorname{tr}\rho\log\rho$.

The book states the assumption of determinism and reversibility as: ``specifying the state of
the system at a particular time is equivalent to specifying the state at a future (determinism)
or past (reversibility) time''~\cite[Ch.~\emph{Classical mechanics}, Sec.~\emph{Reversing
Hamiltonian mechanics}]{book}. In the quantum case, it applies this to pure states, where ``there is a
one-to-one correspondence between initial and final state''~\cite[Ch.~\emph{Quantum mechanics},
Sec.~\emph{Schroedinger equation and unitary evolution}]{book}. The states of maximal
information are the pure states. A mixed final state would leave the outcome undetermined, and
it would also carry less information than the initial state. This gives the definition.

\begin{definition}[DR]\label{def:detrev}
  $\Phi$ is \emph{deterministic} if it maps pure states to pure states, and \emph{reversible} if
  it is injective on the pure states.
\end{definition}

So a deterministic $\Phi$ assigns to each initial state one final state, and a reversible one
lets the final state determine the initial one. Surjectivity is not part of the definition; it
follows (Theorem~\ref{thm:detrev}).

\begin{theorem}\label{thm:detrev}
  For $\Phi$ as above the following are equivalent:
  \begin{enumerate}
    \item[(i)] $\Phi$ is deterministic and reversible (DR);
    \item[(ii)] $\Phi$ is a bijection of $\mathcal P(\mathcal H)$;
    \item[(iii)] $\Phi$ is a bijection of $\mathcal D(\mathcal H)$, an isomorphism of the
      state space in the sense of~\cite[Preface, Sec.~\emph{Project overview}]{book};
    \item[(iv)] $S(\Phi(\rho)) = S(\rho)$ for every $\rho\in\mathcal D(\mathcal H)$ (DR-INFO);
    \item[(v)] $\Phi(\rho) = W\rho W^{-1}$ for a unitary or an anti-unitary $W$.
  \end{enumerate}
\end{theorem}

\begin{corollary}\label{cor:detrev}
  For an evolution $\Psi_{t_0\to t_1}$, $t_0\le t_1$, as in the brief's time-dependent setting,
  DR-INFO holds if and only if every $\Psi_{t_0\to t_1}$ is deterministic and reversible.
\end{corollary}

\subsection{Proof}

\begin{lemma}[Anti-unitary conjugation]\label{lem:antiunitary}
  If $W$ is anti-unitary and $\rho$ is a density operator, then $S(W\rho W^{-1}) = S(\rho)$.
\end{lemma}

\begin{proof}
  Let $\rho = \sum_i \mu_i |v_i\rangle\langle v_i|$ with $(v_i)$ an orthonormal eigenbasis and
  $\mu_i\ge0$. Then $W\rho W^{-1} = \sum_i \mu_i |Wv_i\rangle\langle Wv_i|$, as the $\mu_i$ are
  real. The vectors $Wv_i$ form an orthonormal basis, so the unitary $U$ with $Uv_i = Wv_i$ gives
  $W\rho W^{-1} = U\rho U^{-1}$. The entropy depends only on the characteristic polynomial, which
  unitary conjugation preserves.
\end{proof}

\begin{proof}[Proof of Theorem~\ref{thm:detrev}]
  Since $\Phi$ preserves equal mixtures and maps density operators to density operators, it is
  affine on $\mathcal D(\mathcal H)$.

  (iv) $\Rightarrow$ (i). These are the brief's Steps~1 and~2. The pure states are exactly the
  density operators of zero entropy, so $\Phi$ maps them to pure states. The entropy of the equal
  mixture of two pure states $\psi, \phi$ is a strictly decreasing function of the transition
  probability $|\langle\psi,\phi\rangle|^2$, so $\Phi$ preserves transition probabilities. If
  $\Phi(|\psi\rangle\langle\psi|) = \Phi(|\phi\rangle\langle\phi|)$, the images have transition
  probability $1$, hence so do $\psi$ and $\phi$, and $|\psi\rangle\langle\psi| =
  |\phi\rangle\langle\phi|$.

  (i) $\Rightarrow$ (v). This is Kadison's theorem~\cite{kadison}, which assumes a bijection, in
  a version for maps that are only injective on the pure states. An affine map of $\mathcal D(\mathcal H)$ that is injective on pure states
  and maps them to pure states preserves orthogonality, and then all transition probabilities.
  Wigner's theorem gives a unitary or anti-unitary $W$ with $\Phi(|\psi\rangle\langle\psi|) =
  |W\psi\rangle\langle W\psi|$. As $\Phi$ is affine and every density operator is a convex
  combination of pure states, $\Phi(\rho) = W\rho W^{-1}$.

  (v) $\Rightarrow$ (ii) $\Leftrightarrow$ (iii). Conjugation by $W$ is a bijection of
  $\mathcal D(\mathcal H)$ with inverse conjugation by $W^{-1}$. An affine bijection of
  $\mathcal D(\mathcal H)$ maps its extreme points, the pure states, bijectively onto themselves.
  Conversely a bijection of the pure states satisfies (i), hence (v), hence (iii).

  (ii) $\Rightarrow$ (i) is immediate, and (v) $\Rightarrow$ (iv) is unitary invariance together
  with Lemma~\ref{lem:antiunitary}.
\end{proof}

\begin{proof}[Proof of Corollary~\ref{cor:detrev}]
  Apply Theorem~\ref{thm:detrev}, (i) $\Leftrightarrow$ (iv), to each $\Psi_{t_0\to t_1}$.
\end{proof}

\begin{remark}
  For a single map the anti-unitary case in (v) is genuine: complex conjugation in a basis is
  deterministic, reversible and preserves the entropy. Along a continuous evolution with
  $\Psi_{t\to t} = \mathrm{id}$ and the composition law it is excluded, and every
  $\Psi_{t_0\to t_1}$ is unitary. As in the brief's proof of Theorem~2, this step needs
  continuity in time.
\end{remark}

\begin{remark}
  In dimension at least two neither half of Definition~\ref{def:detrev} implies DR-INFO alone.
  The constant map onto one pure state is deterministic but not reversible, and it lowers the
  entropy of every mixed state. The depolarizing map $\rho\mapsto\frac12\rho +
  \frac1{2n}\mathbb 1$ is injective on all of $\mathcal D(\mathcal H)$ but not deterministic,
  and it raises the entropy of every pure state. These examples are not formalized.
\end{remark}

\begin{remark}
  Finite dimension is used in (i) $\Rightarrow$ (ii): in infinite dimension an isometry that is
  not unitary is deterministic and reversible in the sense of Definition~\ref{def:detrev} and
  preserves the entropy, but it is not a bijection of the pure states. The companion brief
  \emph{Multiple equivalent characterizations of the Schr\"odinger equation} formalizes such an
  example, the semigroup of right shifts on $L^2(0, \infty)$. Whether (i) $\Leftrightarrow$
  (iv) holds in infinite dimension is not addressed here.
\end{remark}

\subsection{Lean}

Module \lean{UnitaryEntropy.FurtherResults.DeterminismReversibility}, namespace
\lean{UnitaryEntropy.FurtherResults}. Pure states and Kadison's theorem come from
\lean{ForPhyslib/QuantumMechanics/Wigner/}, and Steps~1 and~2 from the brief's main file.

\begin{itemize}
  \item \lean{IsDetRev}, \lean{PreservesEntropy}: Definition~\ref{def:detrev} and DR-INFO.
    The pure states are \lean{QuantumMechanics.Wigner.IsPure}, (ii) is
    \lean{QuantumMechanics.Wigner.BijOnPure} and (v) is
    \lean{QuantumMechanics.Wigner.ImplementedByUnitaryOrAntiunitary}.
  \item \lean{exists_conjOp_eq_conjOp}, \lean{vonNeumann_conjOp_antiunitary}:
    Lemma~\ref{lem:antiunitary}; \lean{vonNeumann_conjOp}: (v) $\Rightarrow$ (iv).
  \item \lean{isDetRev_of_preservesEntropy}: (iv) $\Rightarrow$ (i), from
    \lean{UnitaryEntropy.pure_and_born_of_preserves}.
  \item \lean{QuantumMechanics.Wigner.AffineOnDensities.implemented_of_injOn}: (i)
    $\Rightarrow$ (v); \lean{QuantumMechanics.Wigner.kadison_tfae}: (ii) $\Leftrightarrow$
    (iii) $\Leftrightarrow$ (v).
  \item \lean{detRev_tfae}, \lean{isDetRev_iff_preservesEntropy}: Theorem~\ref{thm:detrev}.
  \item \lean{detRev_iff_drInfo}: Corollary~\ref{cor:detrev}, for
    \lean{SchroedingerEvolution.TwoParameter.IsEvolution} and the brief's
    \lean{UnitaryEntropy.TwoParameter.DrInfo}.
\end{itemize}

\begin{thebibliography}{9}
\bibitem{brief} G.~Carcassi, \emph{Schr\"odinger equation is equivalent to information
  conservation}, technical brief, Assumptions of Physics,
  \url{https://assumptionsofphysics.org/resources/briefs}.
\bibitem{book} G.~Carcassi and C.~A.~Aidala, \emph{Assumptions of Physics},
  \url{https://assumptionsofphysics.org}.
\bibitem{kadison} R.~V.~Kadison, \emph{Transformations of states in operator theory and
  dynamics}, Topology~3, suppl.~2 (1965), 177--198.
\end{thebibliography}

\end{document}
